\section{Exact representation and completion of the theorem}
\label{sec:exact-completion}

We now pass from the positive finite Fourier coefficient to an exact rational identity, and then complete the reductions announced in \cref{sec:introduction}.

\subsection{The avoiding unit representation}

Let $W$ be the Bernoulli weight of subsets $A\subset\mathcal E$ satisfying
\[
 \sum_{e\in A}\frac{L}{e}\equiv\frac{L}{b}\pmod L.
\]
By \cref{eq:reciprocal-fourier,thm:positive-fourier-numerator},
\[
 LW=\sum_{h\bmod L}F(h)
\]
has positive real part.  The left side is real and nonnegative, so $W>0$.  Hence there exists a subset $A\subset\mathcal E$ with
\begin{equation}
\label{eq:target-congruence}
 \sum_{e\in A}\frac{L}{e}\equiv\frac{L}{b}\pmod L.
\end{equation}
Dividing by $L$, the difference
\[
 d_A=\sum_{e\in A}\frac1e-\frac1b
\]
is an integer.

The exactness step is now deterministic.  By \cref{eq:exact-centering-load},
\[
 0\leq\sum_{e\in A}\frac1e\leq\Lambda<1,
 \qquad
 0<1/b<1.
\]
Therefore
\[
 -1<d_A<1.
\]
Since $d_A\in\mathbb Z$, it follows that $d_A=0$.  Thus
\begin{equation}
\label{eq:avoiding-unit-identity}
 \sum_{e\in A}\frac1e=\frac1b.
\end{equation}
By \cref{prop:arithmetic-capacity}, the elements of $A$ are distinct squarefree semiprimes and avoid the prescribed finite set $\mathcal T$.

We have proved the fixed-denominator statement on which the elementary closure rests.

\begin{theorem}[Avoiding unit representation]
\label{thm:avoiding-unit}
Let $b\geq3$ be squarefree and let $\mathcal T$ be a finite set of positive integers.  There is a finite set $A$ of pairwise distinct squarefree semiprimes, disjoint from $\mathcal T$, such that
\[
 \sum_{e\in A}\frac1e=\frac1b.
\]
\end{theorem}

This is the unique point at which quotient positivity and ambient equality meet.  The Fourier argument proves the congruence \cref{eq:target-congruence}; only afterwards does the strict interval $[0,\Lambda]\subset[0,1)$ exclude every nonzero alias.

\subsection{Numerators and small denominators}

\begin{lemma}[Numerator induction]
\label{lem:numerator-induction}
Suppose an avoiding representation of $1/b$ exists for every finite forbidden set.  Then, for every integer $a\geq1$ and every finite forbidden set $\mathcal T$, there is an avoiding representation of $a/b$ by distinct squarefree semiprimes.
\end{lemma}

\begin{proof}
Induct on $a$.  The case $a=1$ is the hypothesis.  Suppose a finite set $A$ represents $a/b$ and avoids $\mathcal T$.  Apply the avoiding unit statement with forbidden set $\mathcal T\cup A$ to obtain a representation $A'$ of $1/b$ disjoint from $A$.  Then $A\cup A'$ consists of distinct squarefree semiprimes, avoids $\mathcal T$, and has reciprocal sum $(a+1)/b$.
\end{proof}

\begin{lemma}[Denominators $1$ and $2$]
\label{lem:small-denominators}
It is enough to prove \cref{thm:avoiding-unit} for squarefree $b\geq3$.
\end{lemma}

\begin{proof}
For $b=2$, construct successively disjoint avoiding representations of $1/3$ and $1/6$, and use
\[
 \frac12=\frac13+\frac16.
\]
For $b=1$, construct three successive avoiding representations of $1/3$ and add them; equivalently, apply \cref{lem:numerator-induction} with numerator and denominator equal to $3$.  At each stage the previously used denominators are added to the forbidden set.  Once an avoiding unit representation for $1$ or $1/2$ is obtained, \cref{lem:numerator-induction} supplies every positive numerator.
\end{proof}

\subsection{Necessity of squarefreeness}

\begin{lemma}[Squarefree denominator obstruction]
\label{lem:squarefree-necessity}
Let $A$ be a finite set of squarefree positive integers.  Then the reduced denominator of
\[
 \sum_{n\in A}\frac1n
\]
is squarefree.
\end{lemma}

\begin{proof}
Let $M=\operatorname{lcm}(A)$.  Then
\[
 \sum_{n\in A}\frac1n
 =\frac{\sum_{n\in A}M/n}{M}.
\]
The reduced denominator divides $M$.  Since the least common multiple of squarefree integers is squarefree, every divisor of $M$ is squarefree.
\end{proof}

\begin{proof}[Proof of \cref{thm:main}]
If $q$ is a sum of reciprocals of distinct squarefree semiprimes, its reduced denominator is squarefree by \cref{lem:squarefree-necessity}.

Conversely, write $q=a/b>0$ in lowest terms and suppose $b$ is squarefree.  For $b\geq3$, \cref{thm:avoiding-unit,lem:numerator-induction} give a representation of $a/b$ by distinct squarefree semiprimes.  The cases $b=1,2$ follow from \cref{lem:small-denominators}.  This proves the characterization.
\end{proof}

\subsection{Analytic input}

The prime number theorem is the only external analytic result used in the construction.  It supplies prime counts in fixed-ratio intervals and, by Abel summation, the reciprocal load of the prime blocks.  The synchronization, fibre compression, target observability, five-sector decomposition, positive major estimate, no-wrap step and elementary induction are all proved within the article.
